\section{Conceptual Framework}
\label{sec:conceptual-framework}

Our premise is that the evolution of mathematical practice is driven by a continuous feedback loop between two dimensions: \emph{inherited mathematical structure} and \emph{human organizational process}~\cite{avigad2024,commelin2023}. Inherited structure is the historical residue of past organizational labor, stabilized and sedimented under rigorous logical constraints. Conversely, the organizational process is the active, contemporary reconfiguration of this inherited structure, mediated by evolving tooling environments. Through this intergenerational cycle, the dynamic organizational praxis of one era crystallizes into the foundational inheritance of the next, strictly establishing the epistemic boundaries and starting conditions for future mathematical inquiry.

    \textbf{Inherited mathematical structure} constitutes the epistemic residue of successive generations of mathematical labor. Through historical accumulation and rigorous selection, the mathematical community transmits to contemporary practitioners a highly optimized conceptual ecosystem---definitions, theorems, canonical proof pathways, and disciplinary boundaries---that has already undergone extensive logical testing and historical filtering. Externally embodied in formal literature, standardized curricula, and computational libraries, this system is continuously internalized as domain-specific intuition through education and practice.

    \textbf{Human organizational process} constitutes the dynamic, synchronic reconfiguration of this inherited structure. Enacted by human practitioners and, increasingly, computational agents, this labor is strictly mediated by the sociotechnical tooling environment of a given historical epoch. It encompasses both the systematic refactoring of the extant corpus---via conceptual partitioning, modular packaging, and abstraction---and the generative synthesis of novel mathematics, including new definitions, theorems, and proof architectures. This organizational labor, in turn, recalibrates the mathematical intuition and structural baseline inherited by the next generation.

Each revolution in mathematical tooling---from writing to print, from print to computation---has reshaped the balance between these two dimensions. New tools do not make prior mathematics irrelevant; rather, their utility is bounded by the mathematical framework one brings to them~\cite{demoura2026proofassistants}. The central insight of this analytical framework is that \module{Mathlib} functions as a \emph{transitional artifact}: a large-scale systematic migration of traditional mathematical knowledge---historically bound by the analog constraints of print media and physical locality---into a digital sociotechnical regime characterized by formal verification, decentralized collaboration, and version control. Consequently, its architecture simultaneously embeds the structural signatures of both epochs. Crucially, the proof assistant that mediates this migration is itself a programming language: Lean is simultaneously a proof checker and a general-purpose language whose components mutually reinforce one another. This is causally upstream of graph topology---because Lean is programmable, typeclass instances and coercions become first-class graph nodes, producing the language-infrastructure layer that dominates in-degree (\S\ref{sec:thm-types}). The typeclass mechanism, in particular, represents a fundamental tool-mediated compromise: mathematical knowledge that exists as implicit shared understanding among human practitioners (e.g., ``the integers form a ring'') must be encoded as explicit \emph{instance declarations} to enable automated reasoning, creating an entire stratum of nodes in $\mathcal{D}$ that has no traditional-mathematical counterpart (\S\ref{sec:typeclass-synth}). More broadly, Lean's type system introduces structural layers with no traditional-mathematical counterpart: coercion chains enable automatic type-casting (e.g., $\mathbb{N} \hookrightarrow \mathbb{Z} \hookrightarrow \mathbb{Q}$), the \texttt{extends} mechanism builds structure inheritance hierarchies that organize the algebraic tower, and \texttt{deriving} handlers auto-generate typeclass instances without human intervention. Each mechanism deposits its own signature in the dependency graph.

At the level of logical structure (premise dependencies), the network encodes accumulated mathematical inheritance, yet this foundation has been fundamentally reshaped by the benefits of contemporary tools, specifically Lean's dependent type theory and tactic engine. This migration requires a strict distinction between semantic content and syntactic form. The inherited structure dictates the content: the canonical theorems and macroscopic dependencies remain invariant. Conversely, the contemporary medium dictates the form: microscopic dependency chains, proof architectures, and type-theoretic representations are entirely reconfigured to satisfy the compiler. Therefore, the premise dependency network is neither an isomorphic transcription of traditional mathematics nor a creation \textit{ab initio}, but rather a deeply \emph{tool-mediated migration product}. The compiler further attaches internal metadata to each declaration---definitional heights governing reduction behavior, tactic profiles recording proof strategies---providing fine-grained signals that characterize the tool's imprint on every formalized result. Our metrics, accordingly, are a snapshot of a particular production regime, not an eternal structural law; the temporal contingency of our measurements is acknowledged in \S\ref{sec:limitations}.

Each raw dependency graph, however, is not a pure embodiment of one dimension but an \emph{entanglement} of both inherited structure and organizational process. The graph decompositions defined in \S\ref{sec:definitions} serve as systematic disentangling tools---projections that isolate one dimension from the superposition. At the declaration level, the statement/proof decomposition ($E_S \sqcup E_P \sqcup E_{SP}$) separates what a theorem \emph{asserts} (statement dependencies, largely invariant across proof strategies---inherited structure) from how it is \emph{proved} (proof dependencies, shaped by tactics, conventions, and tooling---organizational process). The explicit/synthesized partition ($E_{\mathrm{explicit}} \sqcup E_{\mathrm{synth}}$) provides the most direct operationalization: edges the mathematician wrote versus edges the elaborator inserted. The typeclass instance graph $G_{\mathrm{tc}}$ isolates the pure tool-infrastructure layer---a form of organization that has no counterpart in traditional mathematics. Finer projections of this layer---the coercion graph $G_{\mathrm{coe}}$, the structure inheritance graph $G_{\mathrm{ext}}$, and auto-derived instance edges $E_{\mathrm{auto}}$---decompose it further by mechanism, distinguishing implicit type-casting from algebraic hierarchy from compiler-generated boilerplate. At the module level, the visibility graph $G_{\mathrm{mod}}^{\mathrm{pub}}$ separates intentional API design from all transitive dependencies, the build graph $G_{\mathrm{build}}$ captures the pure engineering dimension of compilation cost, and import utilization measures how efficiently each module consumes what it imports---revealing whether module boundaries are drawn at natural semantic joints or at arbitrary organizational conventions. Finally, declaration attributes ($\alpha$) enrich nodes with metadata---proof complexity, definitional height, universe polymorphism, tactic usage profiles---that spans both dimensions. Collectively, these decompositions convert the binary inherited/organizational distinction from an interpretive lens into a measurable, falsifiable framework.

Similarly, the library's macro-organizational structure (module partitioning and namespace taxonomies) reflects the computational constraints of modern software development, yet it remains heavily influenced by historical inertia. Top-level directories such as \module{Mathlib.Analysis}, \module{Mathlib.Topology}, and \module{Mathlib.Algebra} map directly onto classical disciplinary boundaries, demonstrating that legacy cognitive classifications persist across the medium shift. Concurrently, the demands of the proof assistant have catalyzed the emergence of entirely novel, medium-specific categories---such as \module{Mathlib.Tactic} and \module{Mathlib.Lean}---which have no analogue in traditional mathematics. The gap between what is formalizable in principle and what is feasible in practice is precisely where proof assistant design lives. Our 17.5\% import redundancy, 0.107 cohesion, and 153-layer depth are quantifications of this gap: the structural distance between what is logically necessary and what is practically maintained.

The structural friction generated by this transition---manifesting as an incomplete alignment between inherited logical dependencies and contemporary organizational boundaries---provides the primary empirical signal for our research. Tightly interdependent declarations are routinely distributed across distinct modules for compilation efficiency, while logically orthogonal content is often aggregated under a single namespace due to legacy naming conventions. These quantifiable misalignments are the observable traces of the collision between mathematical inheritance and organizational practice. Moreover, the inherited-structure/organizational-process loop is self-reinforcing: compact, well-structured proofs are better training data for AI, and AI-generated proofs in turn reshape the library---a flywheel driven by the proof assistant. Three independent teams---AlphaProof, Aristotle, SEED Prover---converging on Lean/\module{Mathlib} is empirical evidence of this flywheel in action (\S\ref{sec:interplay}). As AI systems generate an increasing share of proofs, our structural metrics become a \emph{trust dashboard}: the human contract resides in the specifications, not the proofs, reframing our measurement program as infrastructure for trust in an AI-augmented library (\S\ref{sec:practice}).

Consequently, \module{Mathlib} should be analyzed as a dual ontology: it is simultaneously an \textit{epistemic object} (whose premise network codifies the logical structure of formalized mathematics) and a \textit{sociotechnical artifact} (whose modular hierarchy codifies the pragmatic labor and constraints of a contemporary community). The graph decompositions make this duality operationally precise---each decomposition is a projection that isolates one dimension from the entanglement. The empirical quantification of the separation and interplay between these two dimensions constitutes the core problem of this paper.
